%! Rational Functions and Holes — Unit 3, Lesson 3.4 — Slides (Beamer)
\documentclass[aspectratio=169, 11pt]{beamer}
\usepackage{precalculus-beamer}   % provides \navyheader{} and \sectionlabel[color]{}

\begin{document}

% TITLE SLIDE
{
\setbeamercolor{background canvas}{bg=navy}
\begin{frame}[plain]
  \vspace{2.8cm}\hspace{0.55cm}%
  \begin{minipage}{0.9\textwidth}
    {\color{goldacc}\bfseries\small LESSON 1.10}\\[0.5cm]
    {\color{white}\bfseries\LARGE Rational Functions and Holes}\\[1.2cm]
    {\color{skymid}\small Precalculus~$\cdot$~Unit 3: Rational Functions}
  \end{minipage}
\end{frame}
}

% HOOK SLIDE
\begin{frame}[t, plain]
  \navyheader{Hook: The GPS Signal Gap}
  \sectionlabel[goldacc]{Think About It}

  Suppose your GPS loses signal for exactly one moment --- at $t = 5$ seconds --- but your
  position function is defined everywhere else.

  \bigskip
  \begin{columns}[T]
    \column{0.48\textwidth}
      \begin{block}{What we know}
        \begin{itemize}\setlength{\itemsep}{6pt}
          \item The output \textit{at} $t = 5$ is missing (undefined).
          \item Outputs for $t$ near 5 are all well-defined.
          \item Those nearby outputs cluster around one value $L$.
        \end{itemize}
      \end{block}

    \column{0.48\textwidth}
      \begin{block}{What this means}
        \begin{itemize}\setlength{\itemsep}{6pt}
          \item We can predict the ``missing'' value: $L$.
          \item This is exactly a \textbf{hole} in the graph.
          \item We write: $\displaystyle\lim_{t \to 5} r(t) = L$
        \end{itemize}
      \end{block}
  \end{columns}

  \bigskip
  \textit{Question:} How is this different from a vertical asymptote?
\end{frame}

% HOLES VS VAs SLIDE
\begin{frame}[t, plain]
  \navyheader{Holes vs.\ Vertical Asymptotes}
  \sectionlabel[goldacc]{Decision Chart}

  Let $r(x)$ have a shared zero at $x = a$. Let $m$ = numerator multiplicity, $n$ = denominator multiplicity.

  \bigskip
  \begin{columns}[T]
    \column{0.48\textwidth}
      \begin{block}{HOLE at $x = a$ \quad (if $m \geq n$)}
        \begin{itemize}\setlength{\itemsep}{5pt}
          \item Common factor cancels.
          \item Graph has a missing point.
          \item $\lim_{x \to a} r(x) = L$ exists and is finite.
          \item Example: $\dfrac{(x-2)(x+3)}{x-2}$ \\ hole at $x = 2$
        \end{itemize}
      \end{block}

    \column{0.48\textwidth}
      \begin{block}{VA at $x = a$ \quad (if $m < n$)}
        \begin{itemize}\setlength{\itemsep}{5pt}
          \item Denominator ``dominates'' near $x = a$.
          \item Outputs grow without bound.
          \item $\lim_{x \to a} r(x) = \pm\infty$.
          \item Example: $\dfrac{x-2}{(x-2)^2}$ \\ VA at $x = 2$
        \end{itemize}
      \end{block}
  \end{columns}
\end{frame}

% FINDING HOLE COORDINATES SLIDE
\begin{frame}[t, plain]
  \navyheader{Finding the Coordinates of a Hole}
  \sectionlabel[goldacc]{4-Step Procedure}

  \begin{columns}[T]
    \column{0.44\textwidth}
      \begin{block}{Steps}
        \begin{enumerate}\setlength{\itemsep}{4pt}
          \item \textbf{Factor} numerator and denominator.
          \item \textbf{Identify} shared zeros and multiplicities.
          \item \textbf{Cancel} the common factor.
          \item \textbf{Substitute} $x = c$ into simplified form to get $L$.
        \end{enumerate}
      \end{block}

    \column{0.52\textwidth}
      \begin{block}{Example 1: $r(x) = \dfrac{x^2 - 4}{x - 2}$}
        \begin{enumerate}\setlength{\itemsep}{3pt}
          \item Factor: $\dfrac{(x-2)(x+2)}{x-2}$
          \item Shared zero $x=2$; $m = n = 1$, so \textbf{hole}.
          \item Cancel $(x-2)$: simplified $= x+2$
          \item $L = 2 + 2 = 4$
        \end{enumerate}
        \medskip
        \textbf{Hole at $(2, 4)$} \quad and \quad $\displaystyle\lim_{x \to 2} r(x) = 4$
      \end{block}
  \end{columns}
\end{frame}

% PARTNER PRACTICE SLIDE
\begin{frame}[t, plain]
  \navyheader{Partner Practice}
  \sectionlabel[goldacc]{Work with a Partner}

  Analyze $\displaystyle f_2(x) = \frac{x^2 - 1}{(x-1)(x+4)}$.

  \bigskip
  \begin{block}{Your Tasks}
    \begin{enumerate}\setlength{\itemsep}{6pt}
      \item Factor the numerator and write the full factored form.
      \item Identify any shared zeros. For each, compare $m$ and $n$ and classify as hole or VA.
      \item Find the coordinates of any holes (cancel, then substitute).
      \item Identify any additional vertical asymptotes.
      \item Write the limit statement for any holes.
    \end{enumerate}
  \end{block}

  \bigskip
  \textit{Be ready to share: coordinates of the hole, the VA equation, and the limit notation.}
\end{frame}

\end{document}
